\documentclass[book.tex]{subfiles}
\begin{document}
\chapter{Combining Partial Differential Equation Solvers with Graph Neural Networks}

\label{chap:pde_gnn}
\noindent\rule{11cm}{0.4pt} \\
\textbf{Prerequisites:} \autoref{chap:qft_basics}  \\
\textbf{Difficulty Level:} ** \\
\noindent\rule{11cm}{0.4pt}
\vspace{5mm}
Instead of embedding a machine learned approximation into a classical PDE solver, we can instead consider the inverse approach where a classical PDE solver is embedded into the message passing updates of a graph neural network  \cite{belbute2020combining}. The graph neural network is run directly on a non-uniform mesh for the fluid problem, while the classical PDE solver operates on a coarse-grained representation of the underlying mesh.

\begin{figure}
    \centering
    % \includegraphics{figures/Differentiable Physics/fluid_ml/pde_gnn/pde_gnn.png}
    \includegraphics{figures/Differentiable Physics/fluid_ml/pde_gnn/CFD_GCN model.png}
    \caption{Architecture for PDE-GNN System}
    \label{fig:my_label}
\end{figure}

\section{Adjoint Method}

An adjoint method can be used to obtain a gradient while using a classical numerical code to perform the PDE solution. The classical code is run on a downsampeld grid which must be upsampled to merge with the original grid
\begin{align}
    U^{(i)} &= \frac{\sum_{j=1}^k w(n_j) D^{(n_j)}}{\sum_{j=1}^k w(n_j)}
\end{align}
where $w$ is a distance metric, $D$ is the downsampled value, and $U$ is the upsampled value.

\section{Handling Mesh Degeneration}
A subtlety is that mesh updates have to be handled so that the mesh doesn't damage itself during updates. 
%"Solving large complex partial differential equations (PDEs), such as those that arise in computational fluid dynamics (CFD), is a computationally expensive process. This has motivated the use of deep learning approaches to approximate the PDE solutions, yet the simulation results predicted from these approaches typically do not generalize well to truly novel scenarios. In this work, we develop a hybrid (graph) neural network that combines a graph convolutional network with an embedded differentiable fluid dynamics simulator inside the network itself. By combining an actual CFD simulator with the graph network, we show that we can both generalize well to new situations and benefit from the substantial speedup of neural network CFD predictions, while also substantially outperforming the coarse CFD simulation alone."

\end{document}